\documentclass[sigconf,nonacm]{acmart}

\usepackage{amsmath,amssymb,amsthm}
\usepackage{booktabs}
\usepackage{listings}
\usepackage{xcolor}

\newtheorem{theorem}{Theorem}
\newtheorem{corollary}{Corollary}
\newtheorem{definition}{Definition}

\lstset{
  basicstyle=\ttfamily\small,
  keywordstyle=\color{blue},
  commentstyle=\color{gray},
  breaklines=true,
  frame=single,
  numbers=left,
  numberstyle=\tiny\color{gray}
}

\title{Faceted Authority Attenuation: A Product Lattice Model for AI Agent Governance}

\author{Tymofii Pidlisnyi}
\affiliation{%
  \institution{AEOESS}
  \city{Los Angeles}
  \state{California}
  \country{USA}
}
\email{signal@aeoess.com}

\begin{abstract}
As AI agents gain the ability to take real-world actions under delegated human authority,
governance systems must enforce that delegated capabilities can only be attenuated, never
amplified. Current implementations enforce this principle through ad hoc, per-dimension
checks---verifying scope inclusion, spend limits, and delegation depth independently.
We show that these heterogeneous constraints are instances of a single structural
principle: agent authority is an element of a product lattice, and delegation is a monotone
function on that lattice. We formalize this as \emph{faceted authority attenuation} over
seven constraint dimensions (scope, spend, depth, time, reputation, values, reversibility),
prove that the monotonic narrowing invariant holds over the product ordering, and
demonstrate that this formalization yields three concrete benefits unavailable to
scalar or ad hoc approaches: (1) structured, machine-readable denial diagnostics
that identify exactly which constraint dimensions failed and why, (2) cryptographically
signed authorization witnesses that capture an agent's position in the authority lattice
at execution time, and (3) near-miss detection that warns when an agent approaches
constraint boundaries before violation occurs. The formalization is implemented in the
Agent Passport System, an open-source protocol with 1,507 tests across 80 modules,
validated through a cross-protocol composition test with five independent working group
members on production infrastructure.
\end{abstract}

\keywords{AI agent governance, capability attenuation, lattice-based access control,
delegation, authorization, multi-agent systems}

\begin{document}
\maketitle

\section{Introduction}

The deployment of LLM-based agents capable of executing code, making purchases, sending
messages, and coordinating with other agents creates a class of authorization problems
that existing access control frameworks were not designed to handle. Unlike traditional
software principals, LLM agents are non-deterministic, prompt-sensitive, and capable of
generating novel action sequences that were never explicitly programmed.

The principle of least privilege~\cite{saltzer1975} and capability-based
security~\cite{dennis1966} established that delegated capabilities should be attenuatable
but not amplifiable. Miller~\cite{miller2006} formalized this as capability attenuation:
a reference holder cannot confer rights to another that the holder does not possess.

In practice, agent governance systems enforce attenuation through multiple independent
checks: is the requested scope a subset of the delegated scope? Does the requested spend
exceed the limit? Is the delegation depth within bounds? Each check is implemented
separately, producing separate error messages, with no formal relationship between them.

This paper makes the following observation: these heterogeneous checks are not independent
properties. They are projections of a single structural principle onto different constraint
dimensions. Agent authority is an element of a product lattice, and delegation is a monotone
function on that lattice.

This observation is mathematically straightforward---product lattices over ordered sets
are standard constructions. The contribution is not the mathematics itself but the
identification that LLM agent governance constraints form such a lattice, and the
demonstration that recognizing this structure yields concrete operational benefits that
ad hoc implementations miss.

\paragraph{Contributions.}
\begin{enumerate}
\item We identify seven constraint dimensions of LLM agent authority and prove they form
      a product lattice with delegation as a monotone function.
\item We define three protocol artifacts enabled by the lattice structure:
      \emph{AuthorizationWitness}, \emph{ConstraintVector}, and \emph{ConstraintFailure}.
\item We present three case studies where faceted narrowing catches violations that
      single-dimensional or scalar authorization misses.
\item We provide an open-source implementation with 1,507 tests, validated through
      cross-protocol composition with five independent implementations.
\end{enumerate}

\paragraph{Scope and claims.} This paper presents a useful formalization for agent
governance, not a novel mathematical construction. We do not claim that product lattices
are new. We claim that recognizing agent authority \emph{as} a product lattice, and
building protocol artifacts that exploit this structure, yields measurable improvements
in denial diagnostics, forensic verifiability, and proactive safety monitoring.

\section{Background}

\subsection{Capability Attenuation}

Dennis and Van Horn~\cite{dennis1966} introduced capability-based addressing, where
access rights are bound to unforgeable tokens. Miller~\cite{miller2006} refined this
into the capability attenuation principle. The Agent Passport System~\cite{aeoess2026}
applies this to AI agent delegation chains with Ed25519 cryptographic enforcement.
The current formulation specifies attenuation as three independent conditions:
\begin{align}
\text{scope}(d_{i+1}) &\subseteq \text{scope}(d_i) \label{eq:scope} \\
\text{spendLimit}(d_{i+1}) &\leq \text{spendLimit}(d_i) \label{eq:spend} \\
\text{depth}(d_{i+1}) &< \text{depth}(d_i) \label{eq:depth}
\end{align}

\subsection{Lattice-Based Access Control}

Denning~\cite{denning1976} showed that information flow security can be modeled as a
lattice. Bell and LaPadula~\cite{bell1973} formalized this for confidentiality.
Sandhu~\cite{sandhu1993} extended lattice models to role-based access control. Our
authority space is structurally identical to Denning's security lattice, with agent
authority replacing security clearance and delegation replacing information flow. The
key difference: Denning's lattice has two dimensions. Our authority lattice has seven.

\subsection{Abstract Interpretation}

Cousot and Cousot~\cite{cousot1977} introduced abstract interpretation as a framework
for sound program analysis. Gateway constraint checking in our system can be understood
through this lens. The concrete domain is the set of all possible agent actions. The
abstract domain is the authority lattice. The abstraction function $\alpha$ maps an
action to the minimum authority element required. The check
$\alpha(\text{action}) \leq_A \text{auth}(\text{agent})$ is a sound abstract
interpretation: if it passes, the action is within the principal's intended authority.

\section{Faceted Authority Attenuation}

\subsection{Authority Dimensions}

\begin{definition}[Authority Space]
An authority space is a product of $n$ partially ordered sets (lattices):
$A = D_1 \times D_2 \times \cdots \times D_n$
where each $D_k$ is a bounded lattice with meet operation $\sqcap_k$,
top element $\top_k$, and bottom element $\bot_k$.
\end{definition}

For the Agent Passport System, $n = 7$:

\begin{table}[h]
\caption{Authority Dimensions}
\label{tab:dimensions}
\begin{tabular}{@{}llll@{}}
\toprule
Dimension & Lattice & $\top$ & $\bot$ \\
\midrule
Scope & $(\mathcal{P}(S), \subseteq)$ & $S$ & $\emptyset$ \\
Spend & $(\mathbb{R}_{\geq 0}, \leq)$ & $\infty$ & $0$ \\
Depth & $(\mathbb{N}, \leq)$ & max & $0$ \\
Time & $([0,\infty), \leq)$ & $\infty$ & $0$ \\
Reputation & $([0,100], \leq)$ & $100$ & $0$ \\
Values & $(\mathcal{P}(V), \subseteq)$ & $V$ & $\emptyset$ \\
Reversibility & $(\{T,C,I\}, \leq)$ & $I$ & $T$ \\
\bottomrule
\end{tabular}
\end{table}

The reversibility dimension captures action permanence: $T$ (tentative, fully reversible),
$C$ (compensable, reversible with cost), $I$ (irreversible). The ordering $T \leq C \leq I$
means an agent authorized for irreversible actions can perform tentative ones, but not
vice versa.

\begin{definition}[Product Lattice Ordering]
The product ordering $\leq_A$ is defined componentwise:
$a \leq_A b \iff \forall k \in [1,n]: a_k \leq_k b_k$.
Two elements may be incomparable if one exceeds the other in some dimensions but not all.
\end{definition}

\begin{definition}[Delegation as Monotone Function]
A delegation function $\delta: A \to A$ is monotone if:
$\forall a \in A: \delta(a) \leq_A a$.
\end{definition}

\subsection{The Narrowing Theorem}

\begin{theorem}[Faceted Monotonic Narrowing]
For any delegation chain $d_1, d_2, \ldots, d_k$, where $\delta_i$ is the delegation
function producing $d_{i+1}$ from $d_i$:
\[
\text{auth}(d_k) \leq_A \text{auth}(d_{k-1}) \leq_A \cdots \leq_A \text{auth}(d_1)
\]
Authority monotonically narrows along the chain in the product lattice ordering.
\end{theorem}

\begin{proof}
Each $\delta_i$ is monotone by Definition~3. The composition of monotone functions
on a partially ordered set is monotone: if $\delta_i(a) \leq_A a$ and
$\delta_{i+1}(\delta_i(a)) \leq_A \delta_i(a)$, then by transitivity of $\leq_A$,
$\delta_{i+1}(\delta_i(a)) \leq_A a$. By induction on chain length, the transitive
closure of delegation is monotone. \qed
\end{proof}

\begin{corollary}[Independent Facet Narrowing]
Each dimension narrows independently:
$\forall k: \text{auth}_k(d_{i+1}) \leq_k \text{auth}_k(d_i)$.
The original three-condition formulation (Equations~\ref{eq:scope}--\ref{eq:depth})
is a special case with $n=3$.
\end{corollary}

\begin{corollary}[Effective Authority as Meet]
The effective authority of an agent with multiple active delegations is:
$\text{effectiveAuth}(\text{agent}) = \bigsqcap_{d \in \text{delegations}(\text{agent})} \text{auth}(d)$.
An agent cannot claim authority from one delegation to exceed limits imposed by another.
\end{corollary}

\section{Protocol Artifacts}

The lattice structure enables three protocol artifacts that ad hoc implementations
cannot produce.

\subsection{ConstraintVector: Runtime Lattice Evaluation}

A ConstraintVector is the runtime representation of authority evaluation over the product
lattice. Each entry corresponds to one dimension $D_k$:

\begin{lstlisting}[language=Java,caption=ConstraintVector type]
interface ConstraintEvaluation {
  facet: ConstraintFacet    // which dimension
  status: ConstraintStatus  // pass|fail|unknown
  headroom?: number|string  // distance from boundary
  failure?: ConstraintFailure
}
\end{lstlisting}

The \texttt{status} field maps to the lattice comparison:
\texttt{pass} means $\text{auth}_k(\text{action}) \leq_k \text{auth}_k(\text{agent})$;
\texttt{fail} means the action exceeds authority in dimension $k$;
\texttt{unknown} means insufficient evidence to evaluate.

\subsection{AuthorizationWitness: Signed Lattice Position}

An AuthorizationWitness is a cryptographically signed snapshot of the agent's position
in the authority lattice at the moment of execution:

\begin{lstlisting}[language=Java,caption=AuthorizationWitness type]
interface AuthorizationWitness {
  delegationId: string
  agentId: string
  constraints: ConstraintEvaluation[]
  primaryFailure?: ConstraintFacet
  timestamp: string
  gatewaySignature: string // Ed25519
}
\end{lstlisting}

The witness is generated by the enforcement gateway (not self-reported by the agent),
creating a cryptographic record of \emph{where} in the authority lattice the agent was
positioned when the action was evaluated. This record is essential for dispute resolution:
it proves not just that the action was authorized, but which constraints were checked,
what headroom existed, and which dimension was the binding constraint.

\subsection{ConstraintFailure: Structured Denial}

When a constraint check fails, the ConstraintFailure captures dimensional diagnostics:

\begin{lstlisting}[language=Java,caption=ConstraintFailure type]
interface ConstraintFailure {
  facet: ConstraintFacet
  limit: string | number   // authority bound
  actual: string | number  // action requirement
  primaryFailure: boolean
}
\end{lstlisting}

This captures the information needed to reconstruct the lattice comparison: the failure
occurred because $\text{actual}_k >_k \text{limit}_k$ in dimension $k$. The
\texttt{primaryFailure} flag identifies the dimension that would have blocked even if all
others passed, enabling meaningful error messages and targeted remediation.

\section{Comparative Evaluation}

We present three scenarios where faceted narrowing catches violations that scalar or
single-dimensional authorization misses.

\subsection{Case 1: Time-Spend Dual Expiry}

\paragraph{Scenario.} An agent holds a delegation with \$500 spend limit and 24-hour TTL.
At hour 20, the agent has spent \$450 and requests a \$30 action. A scalar system checks
spend (\$450 + \$30 $\leq$ \$500, pass) and time (valid at check time) independently.

\paragraph{Scalar result.} Each check passes independently. The action proceeds. The
delegation expires mid-execution. The receipt is generated against an expired delegation.

\paragraph{Faceted result.} The ConstraintVector evaluates all dimensions simultaneously.
The \texttt{time} facet shows headroom of 4 hours, but the near-miss alert fires because
the headroom-to-execution-time ratio is below threshold. The AuthorizationWitness records
both spend headroom (\$20) and time headroom (4h), enabling the gateway to flag the
action as temporally risky before execution.

\paragraph{What the lattice provides.} The product structure makes the interaction between
dimensions visible. A scalar system cannot express ``within budget but temporally fragile.''
The ConstraintVector captures this as a single object with cross-dimensional risk.

\subsection{Case 2: Reputation-Gated Scope Interaction}

\paragraph{Scenario.} An agent has delegation scope \texttt{[finance:read, finance:write]}
and a reputation score of 45 (Operator tier, max spend \$100). The agent requests a
\texttt{finance:write} action costing \$80.

\paragraph{Scalar result.} Scope check: pass. Spend check: \$80 $\leq$ delegation limit,
pass. The action proceeds without surfacing that the reputation tier caps effective spend
at \$100, and this action consumes 80\% of the tier's budget.

\paragraph{Faceted result.} The ConstraintVector shows: scope pass (full headroom), spend
pass (\$20 headroom against delegation), reputation pass but only \$20 headroom against
tier limit. The binding constraint is reputation, not delegation.

\paragraph{What the lattice provides.} The effective authority is the meet:
$\text{effectiveSpend} = \min(\text{delegationLimit}, \text{tierLimit})$. The scalar
system sees two passing checks. The lattice system identifies which dimension is the
actual bottleneck and by how much.

\subsection{Case 3: Values-Reversibility Interaction}

\paragraph{Scenario.} An agent attests to all seven values floor principles including
F-003 (proportionality). It requests an irreversible action (deleting a database) within
its scope and spend limits.

\paragraph{Scalar result.} Scope: pass. Spend: pass. Values attestation: all principles
attested. The action proceeds. The irreversibility is not considered in relation to the
proportionality principle.

\paragraph{Faceted result.} The ConstraintVector evaluates the \texttt{reversibility}
dimension (irreversible, permitted) and the \texttt{values} dimension (proportionality
flagged for irreversible actions). The cross-dimensional interaction triggers escalation.
An action authorized in every scalar dimension individually still requires review because
two dimensions interact.

\paragraph{What the lattice provides.} The interaction between values and reversibility
is visible only in the product structure. The lattice makes cross-dimensional interactions
expressible as constraint relationships rather than ad hoc rules.

\section{Implementation and Validation}

The model is implemented in the Agent Passport System (APS), an open-source cryptographic
governance protocol. SDK v1.26.0: 1,507 tests, 392 suites, 80 modules (48 core + 32
constitutional governance). Published on npm and PyPI. MCP server with 108 tools.
Apache-2.0 license.

The ProxyGateway evaluates every action request against the full ConstraintVector.
Denials produce ConstraintFailure with dimensional diagnostics. Approvals produce
signed AuthorizationWitness. Near-miss alerting tracks headroom across all constraint
dimensions and generates proactive warnings before any violation occurs.

\paragraph{Cross-protocol validation.} The implementation was validated through a live
cross-protocol composition test (Campaign~7) with five independent working group members.
APS governance artifacts were verified through three independently designed specifications
(QSP-1 encrypted transport, DID Resolution~v1.0, Entity Verification~v1.0) on production
infrastructure. Two independent SDK implementations (TypeScript, Python) from separate
working group members completed the same composition test.

\section{Limitations}

\begin{enumerate}
\item The mathematical construction (product lattices over ordered sets) is standard.
      The contribution is the domain mapping and the protocol artifacts, not the
      mathematics itself.
\item The seven dimensions are not proven to be exhaustive. Additional constraint
      dimensions may be needed for specific deployment contexts.
\item The \texttt{unknown} constraint status extends the model beyond classical
      lattice theory. A full treatment using Belnap's four-valued logic~\cite{belnap1977}
      is deferred to future work.
\item The formalization does not address the semantic gap between authorized
      action categories and their real-world effects. A companion paper on
      compliance-complete failure patterns is in preparation.
\item Evaluation is primarily self-referential. Cross-protocol testing with the
      working group addresses this partially but independent evaluation by
      external researchers would strengthen claims.
\item Near-miss alerting thresholds are operator-configured. The system detects
      proximity to boundaries but does not determine optimal boundaries.
\end{enumerate}

\section{Conclusion}

We formalized agent authority as a product lattice and delegation as monotone attenuation
over seven constraint dimensions. This formalization is not mathematically novel---it
applies standard constructions from lattice theory and abstract interpretation to a
new domain. The value lies in three concrete protocol artifacts that the lattice
structure enables: authorization witnesses that capture an agent's lattice position at
execution time, constraint vectors that make cross-dimensional interactions visible,
and structured denials that identify exactly which dimensions failed and why.

The formalization unifies what was previously a collection of ad hoc checks into a
single, analyzable structure. We believe this is a useful foundation for agent governance
systems that must enforce capability attenuation across heterogeneous constraint
dimensions while providing machine-readable diagnostics and cryptographic evidence
of authorization state.

All code, tests, and specifications are available at
\url{https://github.com/aeoess/agent-passport-system} under Apache-2.0.

\bibliographystyle{ACM-Reference-Format}
\begin{thebibliography}{12}

\bibitem{aeoess2026}
T.~Pidlisnyi. 2026. Monotonic Narrowing for Agent Authority.
Zenodo. DOI:10.5281/zenodo.18749779

\bibitem{belnap1977}
N.~Belnap. 1977. A useful four-valued logic.
In \emph{Modern Uses of Multiple-Valued Logic}. Springer, 5--37.

\bibitem{bell1973}
D.~Bell, L.~LaPadula. 1973. Secure Computer Systems: Mathematical Foundations.
MITRE Technical Report 2547.

\bibitem{cousot1977}
Patrick Cousot and Radhia Cousot. 1977. Abstract Interpretation: A Unified Lattice
Model for Static Analysis of Programs. In \emph{POPL '77}. ACM, 238--252.

\bibitem{denning1976}
Dorothy E. Denning. 1976. A Lattice Model of Secure Information Flow. \emph{Commun. ACM}
19, 5, 236--243.

\bibitem{dennis1966}
Jack B. Dennis and Earl C. Van Horn. 1966. Programming Semantics for Multiprogrammed
Computations. \emph{Commun. ACM} 9, 3, 143--155.

\bibitem{miller2006}
Mark S. Miller. 2006. \emph{Robust Composition: Towards a Unified Approach to Access
Control and Concurrency Control}. Ph.D. Dissertation. Johns Hopkins University.

\bibitem{mishra2026}
Sanjeev Mishra. 2026. Delegated Agent Authorization Protocol.
Internet-Draft draft-mishra-oauth-agent-grants-01. IETF.

\bibitem{nist2023}
NIST. 2023. AI Risk Management Framework (AI RMF 1.0). NIST AI 100-1.

\bibitem{owasp2024}
OWASP. 2024. AI Vulnerability Scoring System (AIVSS).

\bibitem{saltzer1975}
Jerome H. Saltzer and Michael D. Schroeder. 1975. The Protection of Information in
Computer Systems. \emph{Proceedings of the IEEE} 63, 9, 1278--1308.

\bibitem{sandhu1993}
Ravi S. Sandhu. 1993. Lattice-Based Access Control Models. \emph{Computer} 26, 11, 9--19.

\end{thebibliography}

\end{document}
